\subsection{Grouped ambiguity constraints}
\label{app:core-grouped-cubes}

Within the four-state quotient, a negative observation may distinguish
one marked source from other free inputs. Such constraints retain more
information than their separate Boolean edges, and the additional
implications can themselves be needed in later cuts.

Fix a finite variable set $V$. A \emph{marked cube} $z=(l_z,u_z,j_z)$ has
$l_z\subseteq u_z\subseteq V$ and $j_z\in u_z\setminus l_z$.
Its four-state row has $c$ on $l_z$, $b$ outside $u_z$,
$\pv$ at $j_z$, and $\nv$ at the other free coordinates.
A pure $\wedge,\vee$ formula is \emph{safe} on a budget $Z$ when
it omits $\pv$ at every marked row. Endpoint constants are
allowed here; they simplify away in every nonconstant Boolean function.
Let $\mathcal L_Z$ be the Boolean functions of safe formulas, and put
\[
 x\preceq_Z y\quad\Longleftrightarrow\quad
 g(x)\leq g(y)\text{ for every }g\in\mathcal L_Z.
\]
The safe functions form a bounded sublattice, since two safe values
cannot produce $\pv$ under either operation. Finiteness therefore
makes $\mathcal L_Z$ exactly the upsets of $\preceq_Z$: intersect all
safe functions positive at a vertex to obtain its principal upset,
then take unions of these.

Write $M(S)=\{j_z:z\in S\}$ and let $R_S$ be the initial preorder
on $2^V$ generated by inclusion and the individually marked edges
\[
 A\cup\{j_z\}\longrightarrow A
 \qquad(l_z\subseteq A\subseteq u_z\setminus\{j_z\},\ z\in S).
\]
Every safe Boolean function respects these edges: on the corresponding
one-dimensional cube an essential marked coordinate would force
$\pv$.

\begin{lemma}[A grouped cut]
\label{core:lem:grouped-cut}
Suppose $Q$ is a preorder respected by every $S$-safe Boolean function.
If
\[
 x\setminus M(S)\subseteq y\setminus M(S),\qquad
 xQ u_z,\quad l_zQ y\quad(z\in S),
\]
then $x\preceq_S y$.
\end{lemma}
\begin{proof}
A safe separator, positive at $x$ and negative at $y$, would be
nonconstant on every marked cube, by the two $Q$ premises.
At either Boolean gate, a child producing $\pv$ could be
masked only by a stable absorbing endpoint, which would make the
parent Boolean function constant on that cube. Hence both children
are $S$-safe. Follow a child preserving the separator and repeat.
An endpoint constant cannot separate $x,y$. A separating projection
has its coordinate in $x\setminus y$, hence in $M(S)$, so it is
unsafe at a row marking that coordinate. This is a contradiction.
\end{proof}

This gives an exact recursive theory using one common lower-budget
preorder at each stage. For preorders, write $\vee$ for transitive
closure of union and $\wedge$ for intersection. Let $E_z$ be inclusion
with $u_z\to l_z$ adjoined, and define
$xF_My$ by $x\setminus M\subseteq y\setminus M$.

\begin{proposition}[Joined lower-budget theory]
\label{core:prop:grouped-safety-preorder}
The empty-budget preorder is inclusion. For nonempty $S$, put
$B_S=\bigvee_{z\in S}\preceq_{S\setminus\{z\}}$. Then
\[
 \preceq_S=
 F_{M(S)}\wedge\bigwedge_{z\in S}(B_S\vee E_z).
\]
\end{proposition}
\begin{proof}
Every lower-budget implication remains valid on $S$, so $B_S$ does.
For an edge of the right side already in $B_S$, this proves validity.
Every other edge has, for each $z$, a path using the added reversal;
thus $xB_Su_z$ and $l_zB_Sy$. The grouped-cut lemma proves validity.

Conversely, a pair outside the right side either has a separating
unmarked projection, or is separated by an upset of $B_S\vee E_z$
for some $z$. Such an upset is an upset of
$\preceq_{S\setminus\{z\}}$, so it has a formula safe on that
lower budget. Its two cube endpoints agree because $E_z$ is included.
Stable endpoint semantics then makes the same formula safe at $z$.
Thus it separates the pair on the whole budget $S$.
\end{proof}

The derived information in $B_S$ cannot be replaced by initial marked
edges, even if one takes all raw grouped cuts over every subbudget.
Precisely, call the grouped-cut premises with $Q=R_S$ a \emph{raw cut},
and let $Q_Z$ be the preorder generated by $R_Z$ and all such cuts
for $S\subseteq Z$.

\begin{proposition}[Nested cuts are necessary]
\label{core:prop:grouped-cut-nesting}
On five variables, the following four marked rows, each with exactly
two free coordinates, give
$Q_Z\subsetneq\preceq_Z$:
\[
\begin{array}{c|ccccc}
 &A&B&C&D&E\\\hline
 z_1&\pv&c&\nv&c&b\\
 z_2&\nv&\pv&c&b&c\\
 z_3&c&b&\pv&c&\nv\\
 z_4&b&c&\nv&c&\pv
\end{array}
\]
The monotone Boolean function
$f=ABC\vee BCD\vee ACE\vee ADE$ respects $Q_Z$ but has no safe
formula on $Z$, even allowing observer union.
\end{proposition}
\begin{proof}
Products denote Boolean conjunctions. The function $f$ restricts on
the four cubes to $C,A,E,C$, respectively, so it respects $R_Z$.
Two formulas with Boolean function $f$ are
\[
 T_1=A\wedge\bigl((B\wedge C)\vee(E\wedge(C\vee D))\bigr)
       \vee(B\wedge C\wedge D),
 \qquad
 T_2=(B\wedge C\wedge(A\vee D))
       \vee(A\wedge E\wedge(C\vee D)).
\]
Their values on the four rows are
$(\pv,\nv,\nv,\nv)$ and
$(\nv,\pv,\nv,\nv)$.
Consequently $f$ respects every raw cut on a subbudget omitting
$z_1$ or $z_2$.

For all remaining subbudgets use the initial-edge invariants
\[
 d=BD,\qquad h=BC\vee AE\vee CE.
\]
Their restrictions to the four cubes are $(1,0,0,1)$ and
$(C,1,E,C)$, independent of the respective marked coordinates.
Both therefore respect every $R_S$.
A raw cut using $z_1,z_2$ has $BD\,R_S\,y$ and $CE\,R_S\,y$,
forcing $d(y)=h(y)=1$. But
$dh=BCD\vee ABDE\leq f$, so its conclusion cannot have $f(y)=0$.
Thus every raw cut respects $f$, and so does $Q_Z$.

Here is the nested implication. On the subbudget
$\{z_1,z_3,z_4\}$, the raw grouped cut gives $CE\preceq BD$.
Its upper routes are
\[
 CE\to BCDE\to BCD\to ABCD=u_1,\qquad
 CE\subseteq ACDE=u_3,\qquad CE\subseteq BCDE=u_4;
\]
the first uses the $z_4$ marked edge. Its lower routes are
$BD\to BD$, $AD\to ABD\to BD$ using $z_1$, and $BD\to BD$.
The coordinates removed, $C,E$, are both marked in this subbudget.

Now suppose a $Z$-safe formula separates $ABC$ from $BD$.
It is nonconstant on $z_1$ immediately. It is nonconstant on $z_2$
by $CE\preceq BD$. For $z_3,z_4$, the upper routes are
\[
 ABC\to ABCE\to ACE\to ACDE=u_3,\qquad
 ABC\to ABCD\to BCD\to BCDE=u_4,
\]
using marked edges of $z_2,z_1$, respectively; the lower routes
$AD\to ABD\to BD$ and $BD\to BD$ are initial.
The sole unmarked coordinate $D$ is raised, so the first grouped-cut
premise also holds. The descent therefore gives $ABC\preceq_Z BD$,
contrary to $f(ABC)=1$, $f(BD)=0$.

Finally, distributing the two pure operations over observer union
expands any full term into a union of pure components. If the full
Boolean restriction is $f$, every component has that same restriction:
distinct endpoint outputs would union to $\nv$.
Safety of the union forces every component to be safe. The pure
nonexistence just proved therefore excludes observer union as well.
\end{proof}
